% This is samplepaper.tex, a sample chapter demonstrating the
% LLNCS macro package for Springer Computer Science proceedings;
% Version 2.21 of 2022/01/12
%
\documentclass[runningheads]{llncs}
%
\usepackage[T1]{fontenc}
% T1 fonts will be used to generate the final print and online PDFs,
% so please use T1 fonts in your manuscript whenever possible.
% Other font encondings may result in incorrect characters.
%
\usepackage{graphicx}
% Used for displaying a sample figure. If possible, figure files should
% be included in EPS format.
%
% If you use the hyperref package, please uncomment the following two lines
% to display URLs in blue roman font according to Springer's eBook style:
%\usepackage{color}
%\renewcommand\UrlFont{\color{blue}\rmfamily}
%\urlstyle{rm}
%
\begin{document}
%
\title{Neurosymbolic Injection of LTLf Constraints in Deep Reinforcement Learning Policies}
%
%\titlerunning{Abbreviated paper title}
% If the paper title is too long for the running head, you can set
% an abbreviated paper title here
%
\author{First Author\inst{1}\orcidID{0000-1111-2222-3333} \and
Second Author\inst{2,3}\orcidID{1111-2222-3333-4444} \and
Third Author\inst{3}\orcidID{2222--3333-4444-5555}}
%
\authorrunning{F. Author et al.}
% First names are abbreviated in the running head.
% If there are more than two authors, 'et al.' is used.
%
\institute{Princeton University, Princeton NJ 08544, USA \and
Springer Heidelberg, Tiergartenstr. 17, 69121 Heidelberg, Germany
\email{lncs@springer.com}\\
\url{http://www.springer.com/gp/computer-science/lncs} \and
ABC Institute, Rupert-Karls-University Heidelberg, Heidelberg, Germany\\
\email{\{abc,lncs\}@uni-heidelberg.de}}
%
\maketitle              % typeset the header of the contribution
%
\begin{abstract}
The abstract should briefly summarize the contents of the paper in
150--250 words.

\keywords{First keyword  \and Second keyword \and Another keyword.}
\end{abstract}
%
%
%
\section*{Introduction}
\section{Background}
\subsection{Linear Temporal Logic over finite traces}
We evaluate each LTLf formula on finite trajectories with an explicit termination convention. Given a trace of symbols $\tau=(\sigma_1,\ldots,\sigma_T)$, the implementation consumes all symbols in order, then consumes exactly one explicit end symbol, and finally checks whether the resulting DFA state is accepting. This ``explicit-once'' termination rule is shared across dataset serialization, token-to-symbol mapping, DFA rollout, and logic-loss evaluation.

In early prototypes, end handling was duplicated across modules (e.g., missing end tokens in some paths and extra appended ends in others). That inconsistency produced false unsatisfaction on traces that should satisfy simple safety formulas, and disagreement between hard evaluation and differentiable evaluation. We remove this failure mode by using one adapter-owned end-token convention across the full logic stack.

\subsection{Reinforcement Learning}
We study offline sequence RL with transformer policies trained on tokenized transitions. In this setting, logic satisfaction is treated as a first-class objective: we combine supervised next-token prediction with automata-based constraint feedback, and report both task performance and temporal-logic satisfaction metrics.

\section{Related Works}

\section{Method}
\subsection{Implementation Details and Reproducibility}
We standardize experiment execution via a unified command-line interface with smoke-mode validation before full runs. Each run produces structured artifacts, including metrics files, checkpoint snapshots, and diagnostic plots, and these artifacts are logged with run-level metadata for traceability. Dependencies are separated into core logic/modeling, transformer-specific, and benchmark-specific groups to reduce environment conflicts. We use unit tests and CI to validate environment wrappers, trajectory formatting, and automata-based satisfaction evaluation. All runs write standardized configuration and metric artifacts to support aggregation and reproducibility. In our latest server snapshot, the containerized stack ran with Python 3.10.12 and \texttt{torch==2.3.1+cu121}, with CUDA visible on two GPUs.

\subsection{Finite-Trace Logic Stack}
The adapter defines the authoritative token schema and end token identifier used by all components. Dataset generators create full episodes with exactly one end marker, while training windows may or may not contain that marker depending on where the window is sampled. Canonical satisfaction evaluation is defined on complete finite traces: traces without an end marker are treated as incomplete by default, and therefore unsatisfied in canonical mode.

\subsection{Differentiable Satisfaction with Hard and Soft Traces}
We expose two aligned interfaces for satisfaction checking. The hard path maps token IDs to symbol IDs and returns exact Boolean satisfaction per trajectory. The soft path maps token distributions to symbol distributions and returns acceptance probabilities for differentiable training. Let $p_t(k)$ be token probabilities at step $t$ and let $m_t(k)$ map token $k$ to symbol $s$ at position $t$. Then symbol probabilities are aggregated as $P_t(s)=\sum_{k:m_t(k)=s} p_t(k)$ before DeepDFA rollout. This keeps automata transitions differentiable while preserving consistency with exact DFA semantics; one-hot token distributions recover the hard-path result up to numerical tolerance.

\subsection{Automaton-Aware Constrained Decoding}
At evaluation time, we support greedy decoding, unconstrained beam search, and automaton-aware constrained beam search. Each beam maintains an explicit DFA state that is updated after each decoded transition token group. For a partial beam $b_{1:t}$ with DFA state $q_t$, we score
\begin{equation}
\mathrm{Score}(b_{1:t})=\sum_{i=1}^{t}\log p(a_i \mid h_i)+\lambda_{\mathrm{sat}} \cdot \mathbf{1}[q_t \in F],
\end{equation}
where $F$ is the accepting-state set and $\lambda_{\mathrm{sat}}$ controls soft reranking toward beams consistent with the temporal specification. In constrained mode, we additionally apply hard pruning: beams that enter identified rejecting sink states are removed. This hybrid design preserves probabilistic search while enforcing automaton-level safety constraints when available.

\section{Preliminary Experiments}
\subsection{Task Specifications}
We use environment-specific named specification presets to avoid fragile hand-written formulas in experiment scripts. Propositions follow a position-aware naming scheme such as \texttt{s0\_bin22} for state bins and \texttt{a0\_bin3} for action bins. Runtime preset selection is currently enabled for ColourBomb and NRM navigation; FrozenLake and AntMaze presets are compile-tested placeholders in Milestone 1.

\begin{table}[t]
\caption{Specification preset examples used in Milestone 1.}
\label{tab:spec_presets}
\centering
\small
\begin{tabular}{llll}
\hline
Env & Spec name & Formula & Purpose \\
\hline
ColourBomb & avoid\_single\_bomb\_22 & \texttt{G(!(s0\_bin22))} & avoid bomb state \\
NRM nav & avoid\_unsafe & \texttt{G(!(s0\_bin11 | s0\_bin18))} & avoid unsafe cells \\
FrozenLake & placeholder\_safe & \texttt{G(!(s0\_bin5))} & compile-tested placeholder \\
AntMaze & placeholder\_reach & \texttt{F(s0\_bin2)} & compile-tested placeholder \\
\hline
\end{tabular}
\end{table}

\subsection{DFA Inspection and Diagnostics}
To audit specification complexity and debugging behavior, we export per-spec automata summaries and graph representations (\texttt{.dot}, optional \texttt{.png}). These utilities report state counts, accepting states, sink states, and symbol-set size, and are used to sanity-check compiled formulas before full training runs.

\subsection{Metrics}
We report four primary metrics for policy reliability under logic constraints. \textbf{Return} is the mean episodic reward over rollout episodes. \textbf{Violation rate} is reported both at episode level (fraction of episodes whose finalized trace violates the specification) and step level (fraction of unsafe steps under environment-specific unsafe predicates). \textbf{Hard LTLf satisfaction rate} is the fraction of rollout traces accepted by the compiled DFA under explicit finite-trace termination semantics. \textbf{Soft satisfaction score} is the mean acceptance probability from the differentiable automaton (DeepDFA) evaluated on symbol probabilities. Unless stated otherwise, all reported reliability metrics are computed on executed environment rollouts; decoded token prefixes are used only to select actions during planning.

On a fixed-checkpoint NRM navigation smoke comparison (\texttt{16} episodes, spec \texttt{G(!(s0\_bin0))}), greedy decoding produced \texttt{violation\_rate=1.0} and \texttt{satisfaction\_rate=0.0}, while constrained beam decoding with hard reject-sink pruning produced \texttt{violation\_rate=0.0} and \texttt{satisfaction\_rate=1.0}. This run pair illustrates the expected reliability gain from automaton-aware decoding, with a measured runtime tradeoff (about \texttt{0.60}s greedy vs \texttt{16.99}s constrained in this smoke configuration).

\begin{figure}[t]
\centering
\fbox{\parbox{0.9\linewidth}{
\textbf{Placeholder figure.} Replace with exported automaton image, e.g., \texttt{paper/figures/dfa\_cb\_avoid\_single\_bomb\_22.png}.\\
Suggested content: DFA compiled from a reach-while-avoid style finite-trace specification.
}}
\caption{Example DFA compiled from an environment preset specification.}
\label{fig:dfa_inspection}
\end{figure}
\section{Conclusion}
\begin{credits}
\subsubsection{\ackname} A bold run-in heading in small font size at the end of the paper is
used for general acknowledgments, for example: This study was funded
by X (grant number Y).

\subsubsection{\discintname}
It is now necessary to declare any competing interests or to specifically
state that the authors have no competing interests. Please place the
statement with a bold run-in heading in small font size beneath the
(optional) acknowledgments\footnote{If EquinOCS, our proceedings submission
system, is used, then the disclaimer can be provided directly in the system.},
for example: The authors have no competing interests to declare that are
relevant to the content of this article. Or: Author A has received research
grants from Company W. Author B has received a speaker honorarium from
Company X and owns stock in Company Y. Author C is a member of committee Z.
\end{credits}
%
% ---- Bibliography ----
%
% BibTeX users should specify bibliography style 'splncs04'.
% References will then be sorted and formatted in the correct style.
%
% \bibliographystyle{splncs04}
% \bibliography{mybibliography}
%
\begin{thebibliography}{8}
\bibitem{ref_article1}
Author, F.: Article title. Journal \textbf{2}(5), 99--110 (2016)

\bibitem{ref_lncs1}
Author, F., Author, S.: Title of a proceedings paper. In: Editor,
F., Editor, S. (eds.) CONFERENCE 2016, LNCS, vol. 9999, pp. 1--13.
Springer, Heidelberg (2016). \doi{10.10007/1234567890}

\bibitem{ref_book1}
Author, F., Author, S., Author, T.: Book title. 2nd edn. Publisher,
Location (1999)

\bibitem{ref_proc1}
Author, A.-B.: Contribution title. In: 9th International Proceedings
on Proceedings, pp. 1--2. Publisher, Location (2010)

\bibitem{ref_url1}
LNCS Homepage, \url{http://www.springer.com/lncs}, last accessed 2023/10/25
\end{thebibliography}
\end{document}
